\documentclass[11pt]{article}

\usepackage[a4paper,margin=1in]{geometry}
\usepackage{amsmath,amssymb,amsthm,mathtools}
\usepackage{booktabs}
\usepackage{enumitem}
\usepackage{xcolor}
\usepackage[colorlinks=true,linkcolor=blue!45!black,urlcolor=blue!50!black]{hyperref}

\newcommand{\R}{\mathbb R}
\newcommand{\C}{\mathbb C}
\newcommand{\Id}{\mathrm I}
\newcommand{\dd}{\,\mathrm d}
\newcommand{\cE}{\mathcal E}
\newcommand{\cF}{\mathcal F}
\newcommand{\cG}{\mathcal G}
\newcommand{\cH}{\mathcal H}
\newcommand{\cS}{\mathcal S}
\newcommand{\HS}{\mathrm{HS}}
\DeclareMathOperator{\arsinh}{arsinh}

\newtheorem{proposition}{Proposition}[section]
\newtheorem{theorem}[proposition]{Theorem}
\newtheorem{lemma}[proposition]{Lemma}
\newtheorem{corollary}[proposition]{Corollary}
\theoremstyle{definition}
\newtheorem{problem}[proposition]{Open problem}
\newtheorem{remark}[proposition]{Remark}

\title{Stage 4: physical-coupling resummation and the dressed Grassmannian shell}
\author{Working technical note}
\date{21 July 2026}

\begin{document}
\maketitle

\begin{abstract}
The failure of the coefficientwise quadratic estimate QBC$_1$ does not signal
an instability of the physical Schr\"odinger evolution.  We prove this
directly in the radial scalar model.  The exact forward phase can be removed
without changing real-axis flux; the remaining reflection variable evolves
inside the unit disk by a Riccati equation, and its hyperbolic entropy obeys an
exact transfer identity.  The classical one-dimensional quadratic trace
inequality then gives the required compact-energy entropy bound at physical
coupling.  In particular, the two-slab family that makes the quadratic Taylor
coefficient diverge has a uniformly bounded exact Jost entropy.  We also prove
the dual Pl\"ucker identity
\[
 \Gamma_P(J)=\|\mathop{\bigwedge}\nolimits^N(J^{-1}P)\|^{-1},
\]
which identifies the correct $N$-frame for a direct nonlinear attack.  A
\emph{differential} quotient identity for multiplicative cocycles then cancels the
high--high generator exactly and leaves only the charged low column paired
with a dressed Grassmannian regression state.  A
$U(m,m)$ matrix-ball identity supplies the exact full-determinant prototype,
but also exposes the remaining difficulty: its trace cost includes the
uncharged high--high block.  The next live theorem is therefore a
flux-corrected, fixed-$N$ Grassmannian shell estimate for the gauged inverse
frame, not another coefficientwise Born bound.
\end{abstract}

\noindent
\fcolorbox{blue!65!black}{blue!5}{\parbox{0.94\textwidth}{
\textbf{Status.}  The dual Pl\"ucker identity, scalar gauge equations, flux
identity, scalar physical-coupling entropy theorem, uniform control of the
two-slab family, differential high--high cancellation, and the $U(m,m)$
matrix-ball identity are proved below.  The
fixed-$N$ relative shell inequality is formulated precisely but remains open.
No claim about Simon's conjecture is made in this note.
}}

\section{What the failed quadratic test actually says}

For the radial $d=3$ potential
\begin{equation}
 q_\Lambda(r)=\Lambda^{-3/4}\mathbf 1_{[A,A+\Lambda]}(r)
 +b\mathbf 1_{[A+\Lambda+1,A+\Lambda+1+\delta]}(r),
 \label{eq:two-slab-potential}
\end{equation}
the positive energy integral of the second Taylor coefficient of the gauged
Jost logarithm grows like $\Lambda^{1/4}$, while
\begin{equation}
 \int_1^\infty |q_\Lambda(r)|^2\dd r
 =\Lambda^{-1/2}+b^2\delta.
 \label{eq:two-slab-cost}
\end{equation}
This disproves QBC$_1$.  It does not imply that the exact Jost function at
coupling one grows.  Indeed, the long slab produces a forward phase of size
$\Lambda^{1/4}$.  Expanding that phase in the coupling before composing it
with the outer reflector creates the divergent cross coefficient.  Keeping
the phase exponentiated prevents this nonuniform interchange of the limits
$\alpha\to0$ and $\Lambda\to\infty$.

The replacement target must therefore satisfy three rules:
\begin{enumerate}[label=(\roman*),leftmargin=*]
 \item it is stated directly at the physical coupling $\alpha=1$;
 \item it removes the complete forward evolution before estimating reflection;
 \item it propagates a flux or Grassmannian state between shells, so that an
 outer reflection is not treated as an arbitrary independent coefficient.
\end{enumerate}

\section{The exact dual Pl\"ucker frame}

Let $P$ be the orthogonal projection onto the first $N$ coordinates of
$\C^m$, let $Q=\Id-P$, and let $J\in\mathrm{GL}(m,\C)$.  As in the main
programme, put
\begin{equation}
 \Gamma_P(J)
 :=\frac{|\det J|}
 {\|\bigwedge^{m-N}(QJ)\|}.
 \label{eq:Gamma-def}
\end{equation}
Here the exterior norm is the Euclidean norm of all maximal minors of the
displayed rectangular matrix.

\begin{proposition}[Dual Pl\"ucker identity]
\label{prop:dual-plucker}
For every $J\in\mathrm{GL}(m,\C)$,
\begin{equation}
 \frac{\|\bigwedge^{m-N}(QJ)\|}{|\det J|}
 =\|\bigwedge^N(J^{-1}P)\|.
 \label{eq:dual-plucker}
\end{equation}
Equivalently,
\begin{equation}
 \Gamma_P(J)
 =\|\bigwedge^N(J^{-1}P)\|^{-1}
 =\det\!\left(PJ^{-*}J^{-1}P\big|_{P\C^m}\right)^{-1/2}.
 \label{eq:Gamma-inverse-frame}
\end{equation}
\end{proposition}

\begin{proof}
Choose the standard coordinate bases adapted to $P\oplus Q$.  For every
$N$-element row set $T$, Jacobi's complementary-minor identity gives
\begin{equation}
 \det (J^{-1})_{T,P}
 =\pm\frac{\det J_{Q,T^c}}{\det J}.
 \label{eq:jacobi-complement}
\end{equation}
Squaring the moduli and summing over $T$, the Cauchy--Binet formula identifies
the left collection with the coordinates of
$\bigwedge^N(J^{-1}P)$ and the right collection with those of
$\bigwedge^{m-N}(QJ)/\det J$.  This proves~\eqref{eq:dual-plucker}.
The Gram determinant formula for an exterior volume proves
\eqref{eq:Gamma-inverse-frame}.
\end{proof}

\begin{remark}[Why this matters]
The general target is now an honest analytic exterior vector,
\begin{equation}
 \mathfrak g_P(k)=\bigwedge^N\!\left(J(k)^{-1}P\right)
 \in\bigwedge^N\C^m,
 \qquad
 \log\Gamma_P(k)=-\log\|\mathfrak g_P(k)\|.
 \label{eq:dual-exterior-vector}
\end{equation}
It has no distinguished-minor denominator.  For $N=1$ this reduces to
$\Gamma_P=\|J^{-1}e_0\|^{-1}$, but
Proposition~\ref{prop:dual-plucker} works for every fixed $N$.
\end{remark}

\section{Differential quotient flux and exact high--high cancellation}

The dual identity has an infinitesimal counterpart that already passes the
first low/high algebra test.  Let $J(r)\in\mathrm{GL}(m,\C)$ solve a
left-multiplicative cocycle
\begin{equation}
 J'(r)=\mathcal A(r)J(r),
 \qquad
 \mathcal A=
 \begin{pmatrix}
  \mathcal A_{PP}&\mathcal A_{PQ}\\
  \mathcal A_{QP}&\mathcal A_{QQ}
 \end{pmatrix}
 \label{eq:left-cocycle}
\end{equation}
relative to $P\oplus Q$.  Put
\begin{align}
 X&=PJ,&Y&=QJ,&S&=YY^*,
 \label{eq:XY}\\
 C&=XY^*S^{-1},&E&=X-CY,&\Sigma&=EE^*.
 \label{eq:regression-frame}
\end{align}
The matrix $C:Q\C^m\to P\C^m$ is the least-squares Grassmannian regression
coordinate of the low rows against the high rows.  Unlike a chosen minor, it
is globally defined whenever $J$ is invertible, since $Y$ has full row rank.

\begin{proposition}[Exact differential Grassmannian flux]
\label{prop:differential-flux}
With the notation above,
\begin{align}
 EY^*&=0,
 &\Gamma_P(J)^2&=\det\Sigma,
 \label{eq:residual-volume}\\
 \Sigma'
 &=\mathcal B\Sigma+\Sigma\mathcal B^*,
 &\mathcal B&=\mathcal A_{PP}-C\mathcal A_{QP},
 \label{eq:Sigma-evolution}\\
 \frac{\dd}{\dd r}\log\Gamma_P(J(r))
 &=\Re\operatorname{tr}\!\left(
 \mathcal A_{PP}-C\mathcal A_{QP}\right).
 \label{eq:differential-Gamma}
\end{align}
In particular, neither $\mathcal A_{QQ}$ nor $\mathcal A_{PQ}$ appears
directly in the logarithmic increment.  The displayed generator contains only
the charged column $\mathcal A P$; all high--high evolution is absorbed into
the dressed state $C$.
\end{proposition}

\begin{proof}
The definition of $C$ gives $EY^*=0$.  The Schur complement of $S=YY^*$ in
$JJ^*$ is
\begin{equation}
 XX^*-XY^*S^{-1}YX^*=EE^*=\Sigma.
 \label{eq:Schur-Sigma}
\end{equation}
Hence
\begin{equation}
 |\det J|^2=\det S\det\Sigma,
 \end{equation}
which proves~\eqref{eq:residual-volume}.  From~\eqref{eq:left-cocycle},
\begin{align}
 X'&=\mathcal A_{PP}X+\mathcal A_{PQ}Y,
 &Y'&=\mathcal A_{QP}X+\mathcal A_{QQ}Y.
 \end{align}
Substitute $X=CY+E$ into $E'=X'-C'Y-CY'$.  The result has the form
\begin{equation}
 E'=(\mathcal A_{PP}-C\mathcal A_{QP})E+\mathcal K Y
 \label{eq:E-evolution}
\end{equation}
for a matrix $\mathcal K$ whose value is immaterial.  Since $EY^*=0$, the
$\mathcal K Y$ terms vanish in $E'E^*+EE'^*$, proving
\eqref{eq:Sigma-evolution}.  Taking one half of the logarithmic determinant
derivative and using~\eqref{eq:residual-volume} proves
\eqref{eq:differential-Gamma}.
\end{proof}

\begin{corollary}[$2\times2$ low/high model]
\label{cor:two-channel-flux}
Let $m=2$, $N=1$, write the rows of $J$ as $x,y\in\C^{1\times2}$, and set
\begin{equation}
 c=\frac{xy^*}{\|y\|^2}.
 \end{equation}
Then
\begin{equation}
 \Gamma_P(J)=\|x-cy\|,
 \qquad
 \frac{\dd}{\dd r}\log\Gamma_P(J)
 =\Re(\mathcal A_{11}-c\mathcal A_{21}).
 \label{eq:two-channel-flux}
\end{equation}
Thus an arbitrary high--high entry $\mathcal A_{22}$ cancels exactly.  The
remaining analytic task is to control the physical-coupling product
$c\mathcal A_{21}$ after the complete forward gauge, averaged in energy.
\end{corollary}

\begin{remark}[What is and is not proved]
Proposition~\ref{prop:differential-flux} gives the desired quotient
bookkeeping for a multiplicative $m\times m$ cocycle.  The radial
Schr\"odinger equation is second order, so its shell transfer first lives in a
$2m$-dimensional phase space.  One must still identify the corresponding
Jost/Grassmannian chart and prove that elliptic buffered channels can be
eliminated without changing the safe cost.  The proposition supplies the
algebraic model for that reduction; it is not yet the buffered-shell theorem.
\end{remark}

\section{Exact scalar forward gauge and reflected state}

Consider the scalar half-line equation
\begin{equation}
 -f''(r)+q(r)f(r)=k^2f(r),
 \qquad r>1,
 \label{eq:scalar-schrodinger}
\end{equation}
where $q$ is real and supported in $[1,T]$.  With $x=r-1$, write the outgoing
Jost solution as
\begin{equation}
 f(r,k)=a(r,k)e^{ikx}+b(r,k)e^{-ikx},
 \qquad a(T,k)=1,\quad b(T,k)=0,
 \label{eq:ab-representation}
\end{equation}
and impose the usual variation-of-constants constraint.  Then
\begin{align}
 a'&=p\bigl(a+b e^{-2ikx}\bigr),
 &b'&=-p\bigl(ae^{2ikx}+b\bigr),
 &p(r,k)&=\frac{q(r)}{2ik}.
 \label{eq:ab-system}
\end{align}

Define the complete forward phase and the gauged amplitudes by
\begin{equation}
 \theta(r,k)=\int_r^T p(s,k)\dd s,
 \qquad
 \alpha=e^\theta a,
 \qquad
 \beta=e^{-\theta}b.
 \label{eq:scalar-gauge}
\end{equation}

\begin{proposition}[Exact gauged system and flux]
\label{prop:scalar-gauge}
The gauged variables solve
\begin{align}
 \alpha'
 &=p e^{2\theta}e^{-2ikx}\beta,
 &\beta'
 &=-p e^{-2\theta}e^{2ikx}\alpha,
 \label{eq:gauged-ab}
\end{align}
with $(\alpha,\beta)(T)=(1,0)$.  For real $k>0$,
\begin{equation}
 |\alpha(r,k)|^2-|\beta(r,k)|^2
 =|a(r,k)|^2-|b(r,k)|^2=1.
 \label{eq:scalar-flux}
\end{equation}
Consequently $\rho=\beta/\alpha$ lies in the unit disk and satisfies
\begin{align}
 \rho'
 &=-p e^{-2\theta}e^{2ikx}
   -p e^{2\theta}e^{-2ikx}\rho^2,
 &\rho(T,k)&=0,
 \label{eq:rho-riccati}\\
 (\log\alpha)'
 &=p e^{2\theta}e^{-2ikx}\rho.
 \label{eq:log-alpha}
\end{align}
\end{proposition}

\begin{proof}
Differentiate~\eqref{eq:scalar-gauge}, use $\theta'=-p$, and substitute
\eqref{eq:ab-system}; the diagonal terms cancel exactly.  The Wronskian
$W(\overline f,f)$ is constant and equals $2ik$ at $T$, which proves
\eqref{eq:scalar-flux}.  Since $\theta$ is purely imaginary on the real axis,
the gauge does not change the two moduli.  Dividing the two equations in
\eqref{eq:gauged-ab} gives~\eqref{eq:rho-riccati} and
\eqref{eq:log-alpha}.
\end{proof}

At the boundary, the scalar Jost function is
\begin{equation}
 J_q(k)=f(1,k)
 =e^{-\theta(1,k)}\alpha(1,k)
  \bigl(1+e^{2\theta(1,k)}\rho(1,k)\bigr).
 \label{eq:J-rho}
\end{equation}
The flux identity therefore gives the exact bound
\begin{equation}
 \log^+|J_q(k)|
 \leq \log2-\frac12\log\bigl(1-|\rho(1,k)|^2\bigr).
 \label{eq:J-hyperbolic-bound}
\end{equation}
This is the first nonlinear replacement for the failed Taylor coefficient.

\section{Scalar physical-coupling closure}

Let $a_q(k):=a(1,k)$ be the full-line transmission denominator associated
with the compactly supported extension of $q$ by zero to $(-\infty,1]$.
The flux relation gives
\begin{equation}
 |a_q(k)|^2-|b(1,k)|^2=1,
 \qquad
 |J_q(k)|\leq |a_q(k)|+|b(1,k)|\leq2|a_q(k)|.
 \label{eq:J-versus-a}
\end{equation}
The Deift--Killip transmission estimate, a consequence of the
one-dimensional second trace formula, is
\begin{equation}
 \int_\R k^2\log|a_q(k)|\dd k
 \leq\frac{\pi}{8}\int_\R|q(r)|^2\dd r.
 \label{eq:deift-killip}
\end{equation}
Since $\log|a_q|\geq0$, for each
$K=[k_0,k_1]\Subset(0,\infty)$ this gives
\begin{equation}
 \int_K\log|a_q(k)|\dd k
 \leq \frac{\pi}{8k_0^2}\int_\R |q(r)|^2\dd r
 =:C_K\int_\R |q(r)|^2\dd r.
 \label{eq:second-trace-local}
\end{equation}
For compact bounded real potentials this follows from the classical trace
identity after discarding the bound-state contribution, which has the
favourable sign.  Approximation extends the inequality to real $L^2$
potentials.

\begin{theorem}[Exact scalar entropy at coupling one]
\label{thm:scalar-physical}
For every real compactly supported $q\in L^2(1,\infty)$ and every
$K=[k_0,k_1]\Subset(0,\infty)$,
\begin{equation}
 \int_K\log^+|J_q(k)|\dd k
 \leq |K|\log2+C_K\int_1^\infty|q(r)|^2\dd r.
 \label{eq:scalar-physical-entropy}
\end{equation}
The constant is independent of the support radius.
\end{theorem}

\begin{proof}
Combine~\eqref{eq:J-versus-a} with~\eqref{eq:second-trace-local}.
\end{proof}

\begin{remark}
This closes the decoupled radial $m=N=1$ entropy problem nonperturbatively.
It does \emph{not} close the manuscript's $N=1$ problem, where the constant
spherical harmonic is coupled to arbitrarily many angular channels and the
safe cost must exclude the high--high block.
\end{remark}

\begin{theorem}[The QBC$_1$ two-slab family is physically bounded]
\label{thm:two-slab-physical}
Fix $K=[k_0,k_1]\Subset(0,\infty)$, $A>1$, $\delta>0$, and
$0<b<k_0^2/2$.  For the potentials~\eqref{eq:two-slab-potential},
\begin{equation}
 \sup_{\Lambda\geq\Lambda_0}\sup_{k\in K}|J_{q_\Lambda}(k)|<\infty,
 \qquad
 \sup_{\Lambda\geq\Lambda_0}
 \int_K\log^+|J_{q_\Lambda}(k)|\dd k<\infty.
 \label{eq:two-slab-uniform}
\end{equation}
\end{theorem}

\begin{proof}
For a constant potential $c<k_0^2$ on an interval of length $h$, backward
propagation of the physical state $(f,f')^T$ is given by
\begin{equation}
 P_c(h,k)=
 \begin{pmatrix}
  \cos(\kappa h)&-\kappa^{-1}\sin(\kappa h)\\
  \kappa\sin(\kappa h)&\cos(\kappa h)
 \end{pmatrix},
 \qquad \kappa=(k^2-c)^{1/2}.
 \label{eq:constant-transfer}
\end{equation}
For $k\in K$, $c=b$, and $c=\Lambda^{-3/4}$ with $\Lambda$ large,
$\kappa$ remains in a fixed compact subinterval of $(0,\infty)$.  Hence the
norm of~\eqref{eq:constant-transfer} is bounded independently of $h$.
Free propagation has the same property.  The complete two-slab transfer is a
product of a fixed number of such matrices: a fixed free collar, the long weak
slab, the unit gap, and the fixed outer slab.  Its norm is uniform in
$\Lambda$ and $k$.  The outgoing terminal state has uniformly bounded norm,
so its first component at $r=1$, namely $J_{q_\Lambda}(k)$, is uniformly
bounded.  The entropy assertion follows immediately.
\end{proof}

Theorem~\ref{thm:two-slab-physical} is stronger for this family than the
integrated trace bound: it is pointwise in energy.  It rigorously identifies
the QBC$_1$ divergence as a Taylor nonuniformity.

\section{Exact shell composition in the scalar model}

In the forward/backward amplitude basis, the transfer through any real scalar
shell has the $SU(1,1)$ form
\begin{equation}
 M=\begin{pmatrix}u&v\\ \overline v&\overline u\end{pmatrix},
 \qquad |u|^2-|v|^2=1.
 \label{eq:su11}
\end{equation}
If $z_{\rm out}=b_{\rm out}/a_{\rm out}\in\mathbb D$, then
\begin{equation}
 z_{\rm in}
 =\frac{\overline v+\overline u z_{\rm out}}
 {u+v z_{\rm out}}.
 \label{eq:mobius}
\end{equation}
Direct subtraction gives
\begin{equation}
 1-|z_{\rm in}|^2
 =\frac{1-|z_{\rm out}|^2}{|u+v z_{\rm out}|^2}.
 \label{eq:disk-identity}
\end{equation}
Thus, for the hyperbolic flux entropy
\begin{equation}
 \cH(z):=-\log(1-|z|^2),
 \label{eq:hyperbolic-entropy}
\end{equation}
one has the exact dressed increment
\begin{equation}
 \cH(z_{\rm in})-\cH(z_{\rm out})
 =2\log|u+v z_{\rm out}|.
 \label{eq:exact-shell-increment}
\end{equation}

Equation~\eqref{eq:exact-shell-increment} explains why the incoming reflected
state must remain inside the estimate.  Replacing it by $|u|+|v|$, or taking
positive parts before shells are composed, loses the phase correlation that
the trace formula exploits.  Even though
\begin{equation}
 |u+vz|\leq |u|+|v|=\exp(\arsinh|v|),
 \label{eq:bad-majorant}
\end{equation}
the right side is linear in a small reflection amplitude and therefore cannot
in general be summed against a quadratic $L^2$ cost.  The exact logarithm, its
energy dependence, and the flux telescoping are all essential.

\section{The full matrix prototype and its limitation}

Let
\begin{equation}
 \mathbb M=\begin{pmatrix}A&B\\C&D\end{pmatrix}\in U(m,m),
 \qquad
 \mathbb M^*\begin{pmatrix}\Id&0\\0&-\Id\end{pmatrix}\mathbb M
 =\begin{pmatrix}\Id&0\\0&-\Id\end{pmatrix}.
 \label{eq:umm}
\end{equation}
For a matrix reflection state $R$ in the open matrix ball
$\Id-R^*R>0$, define
\begin{equation}
 R'=(C+DR)(A+BR)^{-1}.
 \label{eq:matrix-mobius}
\end{equation}

\begin{proposition}[$U(m,m)$ flux identity]
\label{prop:umm-flux}
Whenever~\eqref{eq:matrix-mobius} is defined,
\begin{equation}
 \Id-R'^*R'
 =(A+BR)^{-*}(\Id-R^*R)(A+BR)^{-1}.
 \label{eq:matrix-ball-identity}
\end{equation}
Consequently
\begin{equation}
 -\log\det(\Id-R'^*R')
 +\log\det(\Id-R^*R)
 =2\log|\det(A+BR)|.
 \label{eq:matrix-hyperbolic-increment}
\end{equation}
\end{proposition}

\begin{proof}
The $U(m,m)$ relations imply
\begin{align}
 &(A+BR)^*(A+BR)-(C+DR)^*(C+DR)
 =\Id-R^*R.
 \label{eq:umm-subtraction}
\end{align}
Multiply on the left by $(A+BR)^{-*}$ and on the right by
$(A+BR)^{-1}$ to obtain~\eqref{eq:matrix-ball-identity}.  Taking determinants
proves~\eqref{eq:matrix-hyperbolic-increment}.
\end{proof}

This is the exact noncommutative analogue of
\eqref{eq:exact-shell-increment}.  Applied to all $m$ channels, however, its
quadratic trace coefficient is the full matrix cost
\begin{equation}
 \operatorname{tr}(M_V^2)=\|M_V\|_{\HS}^2,
 \label{eq:unsafe-full-cost}
\end{equation}
which contains the uncharged high--high block.  The desired Grassmannian
coefficient is only
\begin{equation}
 \operatorname{tr}(PM_V^2P)=\|M_VP\|_{\HS}^2.
 \label{eq:safe-low-column-cost}
\end{equation}
Thus Proposition~\ref{prop:umm-flux} is a prototype, not the answer.  The
fixed-$N$ theorem must be a relative or quotient version built from the dual
frame in Proposition~\ref{prop:dual-plucker}.

\section{The next live theorem}

For a truncated angular model and a radial cutoff $T$, define the analytic
inverse frame
\begin{equation}
 G_{P;R,L,T}(k)=J_{R,L,T}(k)^{-1}P,
 \qquad
 \mathfrak g_{P;R,L,T}(k)=\bigwedge^N G_{P;R,L,T}(k).
 \label{eq:inverse-frame}
\end{equation}
Let $Z_{R,L,T}(k)$ denote the exact relative operator-WKB forward evolution.
On the real axis $Z$ is unitary.  Therefore the gauged exterior frame
\begin{equation}
 \widetilde{\mathfrak g}_{P;R,L,T}(k)
 =\bigwedge^N\!\left(Z_{R,L,T}(k)^{-1}G_{P;R,L,T}(k)\right)
 \label{eq:gauged-inverse-frame}
\end{equation}
has the same boundary norm and
\begin{equation}
 \log\Gamma_{P;R,L,T}(k)
 =-\log\|\widetilde{\mathfrak g}_{P;R,L,T}(k)\|.
 \label{eq:Gamma-gauged-frame}
\end{equation}

The direct physical-coupling Hardy target is:
\begin{problem}[Gauged exterior Hardy theorem $\mathrm{GVH}_N^{\rm phys}$]
\label{prob:gvhn-physical}
For every fixed $N$ and compact $K\Subset(0,\infty)$, choose a complex
neighbourhood $\Omega_K$ and $k_*\in\Omega_K$.  Under a sufficiently small
$X_d$ tail, prove uniformly in $R,L,T$ that
\begin{align}
 \|\widetilde{\mathfrak g}_{P;R,L,T}(k_*)\|
 &\geq c_{K,N}>0,
 \label{eq:gvhn-interior}\\
 \int_{\partial\Omega_K}
 \|\widetilde{\mathfrak g}_{P;R,L,T}(k)\|^2
 \dd\omega_*(k)
 &\leq C_{K,N}\bigl(1+\cE_N(V)\bigr),
 \label{eq:gvhn-boundary}\\
 \cE_N(V)
 &:=\int_R^\infty\|M_{V(r)}P_N\|_{\HS}^2\dd r.
 \label{eq:EN}
\end{align}
Then subharmonicity of
$\log\|\widetilde{\mathfrak g}_{P;R,L,T}\|$ and
\eqref{eq:Gamma-gauged-frame} give PSR$_N$.
\end{problem}

The shell route should prove the same conclusion without expanding in the
coupling.  Divide the tail into buffered shells $\cS_j$ and let
\begin{equation}
 \Delta_j(k)
 :=\log\Gamma_P^{(j)}(k)-\log\Gamma_P^{(j+1)}(k)
 \label{eq:relative-shell-increment}
\end{equation}
be the exact change when shell $j$ is attached to the already dressed outer
tail.  The required local statement is:
\begin{problem}[Dressed Grassmannian shell estimate]
\label{prob:dressed-shell}
Construct a nonnegative flux functional $\cF_j$, depending on the exact
incoming gauged Grassmannian state, such that
\begin{equation}
 \int_K[\Delta_j(k)]_+\dd k
 \leq C_{K,N}\int_{\cS_j^+}
 \|M_{V(r)}P_N\|_{\HS}^2\dd r
 +\cF_j-\cF_{j+1},
 \label{eq:dressed-shell-target}
\end{equation}
uniformly in the angular cutoff and shell radius.  Here $\cS_j^+$ is the
shell together with the fixed impact-parameter buffer.  The functional must
reduce to the hyperbolic reflected flux in
\eqref{eq:hyperbolic-entropy}--\eqref{eq:exact-shell-increment} in the scalar
model and must not charge the high--high block.
\end{problem}

\section{Immediate falsification programme}

Before any Airy estimate is attempted, a candidate for
\eqref{eq:dressed-shell-target} must pass the following tests.
\begin{enumerate}[label=\textbf{Test \arabic*.},leftmargin=*]
 \item \textbf{Two separated slabs.}  The exact forward phase of the long
 slab must remain exponentiated, and the bound must be uniform for
 \eqref{eq:two-slab-potential}.
 \item \textbf{Many weak coherent reflectors.}  No absolute summation of
 individual reflection amplitudes is allowed.  The estimate must use energy
 averaging and flux telescoping, since coherent transfer can be large at one
 energy while the $L^2$ sum rule remains integrated in energy.
 \item \textbf{A $2\times2$ low/high model.}  The multiplicative model passes
 this test by Corollary~\ref{cor:two-channel-flux}: the high--high generator
 cancels exactly.  The phase-space lift must reproduce the same cancellation
 and contain only $\|M_VP\|_{\HS}^2$.
 \item \textbf{Zeros and chart changes.}  The formulation must use the norm
 of the full exterior vector, or a matrix-ball invariant, rather than divide
 by one Grassmannian coordinate.
 \item \textbf{Turning channels last.}  Only after Tests 1--4 survive should
 the proof be made uniform for $\nu\sim kr$ using the Airy buffer.
\end{enumerate}

The first $2\times2$ multiplicative calculation is now complete: its flux is
the residual Gram determinant in Proposition~\ref{prop:differential-flux}.
The next concrete calculation is its $4\times4$ Schr\"odinger phase-space
lift with one retained channel, an arbitrary unitary high-channel forward
cocycle, and small low-column coupling.  Its purpose is to determine whether
the residual determinant descends through the elliptic Schur complement.  If
it does, the following step is an energy-square-function estimate for the
exact nonlinear factor $C\mathcal A_{QP}$.

\section{Conclusion}

The failed QBC$_1$ test has now produced a sharper positive result.  The
physical scalar evolution is controlled, the two-slab obstruction disappears
after exact resummation, the general Pl\"ucker quotient is the reciprocal
volume of a concrete inverse frame, and its multiplicative differential
identity cancels the high--high generator exactly.  The main unresolved step is no longer
``control higher Born terms.''  It is:
\begin{quote}
Find the fixed-$N$ flux on the gauged inverse Grassmannian frame, prove one
buffered-shell inequality with only the low-column $L^2$ cost, and retain its
energy dependence without coefficientwise expansion.
\end{quote}
This is narrower than the former quadratic programme and has an immediate
$4\times4$ phase-space falsification test.

\begin{thebibliography}{9}

\bibitem{DeiftTrubowitz1979}
P.~Deift and E.~Trubowitz,
\emph{Inverse scattering on the line},
Comm. Pure Appl. Math. \textbf{32} (1979), 121--251.

\bibitem{DeiftKillip1999}
P.~Deift and R.~Killip,
\emph{On the absolutely continuous spectrum of one-dimensional Schr\"odinger operators with square summable potentials},
Comm. Math. Phys. \textbf{203} (1999), 341--347.
\url{https://doi.org/10.1007/s002200050615}

\end{thebibliography}

\end{document}
